
There are several different approaches to classify 
cubic surfaces with 27 lines over finite fields $\bbF_q$ in Orbiter.
Classification means to determine the non-equivalent surfaces under the action of the 
collineation group $\PGGL(4,q)$ of $\PG(3,q).$
The approach described in~\cite{BettenKaraoglu19a} 
relies on Schlaefli's notion of 
a double six as a substructure~\cite{Schlaefli58}.
The approach described in~\cite{Karaoglu18} 
utilizes the relation to non-conical six-arcs in a plane.
A third approach is described in~\cite{KaraogluBetten21}.
All three approaches are available in Orbiter.

\bigskip


Table~\ref{tab:classify:cubic:surfaces} 
lists the Orbiter commands to classify cubic surfaces.
\orbitertableofcommands{classify_cubic_surfaces}


\bigskip

Table~\ref{tab:cubic:surface:classified:activity} 
lists the Orbiter commands that can be performed 
once the classification is completed.
\orbitertableofcommands{classification_of_cubic_surfaces_with_double_sixes_activity}

\bigskip



In $\PG(3,4)$, there is only one type of cubic surfaces with 27 lines. 
It is a member of the Hirschfeld family, 
described in~\cite{Hirschfeld64}.
The following Orbiter command can be used 
to construct this surface and to prove its uniqueness 
for $\bbF_4$. 
The following command utilizes the algorithm 
of~\cite{BettenKaraoglu19a} to do so:

\sourcecode{surface_classify_q4}

The \verb'-report' option creates a latex report.
After some redactions, the report contains the following elements.
\orbiteroutput{GRAPHICS/LATEX/report_suface_classify_q4}

\bigskip


Karaoglu~\cite{Karaoglu18} describes a different algorithm, 
based on non-conical six-arcs and Steiner trihedral pairs. 
The command 

\sourcecode{surface_classify_q4_arc_lifting_two_lines}

classifies all cubic surfaces with 27 lines over 
the field $\bbF_{4}$ using the algorithm of Karaoglu.
The result agrees with the previous algorithm. 
In $\PG(3,4),$ the only surface with 27 lines is the Hirschfeld surface.



\bigskip


The classification of cubic surfaces with double sixes 
can be used to perform isomorphism testing and recognition of surfaces. 
The surfaces chosen in the classification algorithm are called canonical surfaces. 
Any surface is isomorphic to exactly one canonical surface.
The canonical surface is often called the fanonical form.
It is possible to efficiently compute a map that takes a given surface to the associated canonical surface.  
Suppose we want to compute the canonical surface of $F_{a,b,c,d}$ with $(a,b,c,d) = (25,5,5,25)$ over $\bbF_{31}.$ 
We can use the following command to create the surface, 
to classify all surfaces with 27 lines over $\bbF_{31},$ and to establish  
an isomorphism between the given surface and the canonical form.


\sourcecode{Family_general_F31_25_5_5_25_recognize}


\bigskip

The following command creates all surfaces in the form 
$F_{abcd}$ over $\bbF_{9}$ and identifies the isomorphism type of each surface.
The result is written to a csv file.

\sourcecode{surface_F_abcd_identify_q9}


\bigskip


The next example illustrates isomorphism testing for smooth cubic surfaces with 27 lines. 
In order to perform isomorphism testing, the classification of cubic surfaces of that order is required.
Let us consider an example over the field $\bbF_{11}.$
Let us see whether the two surfaces of type $F_{abcd}$ 
with $(a,b,c,d) = (2,4,4,2)$ and $(a,b,c,d) = (5,10,7,5)$ over $\bbF_{11}$ are isomorphic.
Both surfaces have 10 Eckardt points. 
However, having the same number of Eckardt points is only a necessary condition for surfaces to be isomorphic. 
It is not a sufficient condition.
To determine whether the surfaces are isomorphic, 
we can use the following command:

\sourcecode{Family_general_F11_iso_test}

The command creates both surfaces, 
and then performs an isomorphism 
test for the two given surfaces. The surfaces are found to be isomorphic. 
A transformation mapping one surface to the other is computed.

\bigskip

Doing isomorphism testing of cubic surfaces is difficult. 
The approach taken here is the following:
Classify all cubic surfaces over the given field and use the data 
obtained during this classification to compute 
the canonical form of each of the given surfaces. 
The surfaces are isomorphic if and only if the respective canonical forms agree. 
By multiplying the group elements that move each of the surfaces 
to their associated canonical form (one element inverted), 
we obtain an isomorphism mapping one to the other. 
Obviously, the computation of one isomorphism is as costly as the classification of all surfaces.
However, once the classification is finished, isomorphisms can be computed in near constant time, 
as the data from the classification allows to determine 
a canonizing group element in almost constant time. 


\bigskip


It is possible to identify cubic surfaces in special forms. 
For instance, the Cayley normal form is parametrized by 4 field elements $(k,l,m,n)$ as
$$
\left(X_0^2+X_1^2+X_2^2+X_3^2\right)X_3 + \left(aX_0+bX_1+cX_2\right)X_3^2+kX_0X_1X_2+dX_0X_1X_3+eX_0X_2X_3+fX_1X_2X_3,
$$
where 
$$
a=l+ \frac{1}{l}, \;
b=m+ \frac{1}{m}, \;
c=n+ \frac{1}{n}, \;
d=lm+ \frac{1}{lm}, \;
e=nl+ \frac{1}{nl}, \;
f=mn+ \frac{1}{mn}.
$$
The following command identifies the Cayley normal form among all cubic surfaces with 27 lines for $q=7.$

\sourcecode{surface_classify_q7_sweep_Cayley}

The command writes two csv files. The first maps $(k,l,m,n)$ to isomorphism types of surfaces. 
The second file writes for each isomorphism type a suitable $(k,l,m,n)$ pair.
There is only one cubic surface over $\bbF_7.$
The computation shows that it can be obtained in the Cayley form with parameters $(k,l,m,n) = (1,2,2,2)$.
This is not the only way. 
In total, there are 64 parameter tuples which all produce isomorphic copies of the same surface.

\bigskip




